% Version A: 141 whitespace-delimited words
\subsection{Sensitivity Analysis of AO Stopping Parameters}
The outer cap $R_{\max}$ limits AO runtime, whereas $\epsilon$ controls termination based on objective change. We tested sixteen combinations of $R_{\max}\in\{4,6,8,10\}$ and $\epsilon\in\{10^{-6},10^{-7},10^{-8},10^{-9}\}$ using 147 compatible exact source pairs at 26 distances. Each setting retained the original coordinate searches, four initializations, and its own previous-distance warm starts. The baseline reproduced all 129 recorded starts and 283 outer iterations. Every setting selected identical parameter pairs and secret-key rates at all distances, matching the finite-catalog maxima. Increasing the cap or tightening the threshold to $10^{-9}$ produced no improvement. However, no coordinate search completed, and all starts requested unavailable source moments. Thus, the observed invariance supports retaining $R_{\max}=6$ and $\epsilon=10^{-8}$ as computational safeguards for this catalog, not as convergence guarantees. The threshold measures objective change rather than absolute SKR accuracy; additional exact sources and completed coordinate refinement are required before claiming continuous numerical optimality.

% Version B (alternative, do not include both): 296 words
\subsection{Sensitivity Analysis of AO Stopping Parameters}
The AO stopping parameters control distinct mechanisms: $R_{\max}$ bounds the number of complete outer cycles, while $\epsilon$ limits the change in the objective between successive cycles. The implementation also uses this objective tolerance for monotonicity warnings and endpoint improvement checks. Consequently, the experiment applied each tolerance consistently through temporary overrides, including explicit arguments for the inner-search default. Parameter tolerance, the four-stage inner refinement cap, boundary expansion rules, physical model, and security calculations were unchanged.

We first reproduced the baseline $R_{\max}=6$, $\epsilon=10^{-8}$ using the 147-source compatible exact catalog. All 129 starts, 283 recorded iterations, candidate histories, terminal states, and selected solutions were reproduced exactly. We then evaluated sixteen combinations of caps in $\{4,6,8,10\}$ and tolerances in $\{10^{-6},10^{-7},10^{-8},10^{-9}\}$. Each configuration used the original four initializations and an independent previous-distance warm-start chain. The reference distance and 22 km were analyzed with their required distance prefixes, followed by the remaining established grid.

Across all 26 distances, every configuration selected the same binary64 parameter pair and SKR as the baseline and matched the finite-catalog maximum. Larger caps and the tighter $10^{-9}$ threshold did not improve the selected objective. Nevertheless, no coordinate search satisfied the original completion conditions. Missing exact source moments blocked each start; therefore, unchanged solutions cannot establish convergence over the continuous parameter domain. Inner searches reconstruct local brackets from the catalog rather than carrying refinement state between outer cycles. Repeating a blocked coordinate search cannot replace unavailable midpoint evaluations, and a zero objective increment can trigger termination before a larger cap becomes active.

These results justify describing $R_{\max}=6$ and $\epsilon=10^{-8}$ as conservative computational safeguards with finite-catalog invariance. They do not establish theoretical optimality, continuous numerical convergence, or a certified absolute SKR error. Cache-only timings and missing-source demand are reported separately because cached lookups do not measure arbitrary-precision source-generation cost.
